%-------------------------
% Resume in Latex
% Author : Jake Gutierrez
% Based off of: https://github.com/sb2nov/resume
% License : MIT
%------------------------

\documentclass[letterpaper,11pt]{article}

\usepackage{latexsym}
\usepackage[empty]{fullpage}
\usepackage{titlesec}
\usepackage{marvosym}
\usepackage[usenames,dvipsnames]{color}
\usepackage{verbatim}
\usepackage{enumitem}
\usepackage[hidelinks]{hyperref}
\usepackage{fancyhdr}
\usepackage[english]{babel}
\usepackage{tabularx}
\usepackage{amsmath}
\input{glyphtounicode}

\pagestyle{fancy}
\fancyhf{} % clear all header and footer fields
\fancyfoot{}
\renewcommand{\headrulewidth}{0pt}
\renewcommand{\footrulewidth}{0pt}

% Adjust margins
\addtolength{\oddsidemargin}{-0.5in}
\addtolength{\evensidemargin}{-0.5in}
\addtolength{\textwidth}{1in}
\addtolength{\topmargin}{-.5in}
\addtolength{\textheight}{1.0in}

\urlstyle{same}

\raggedbottom
\raggedright
\setlength{\tabcolsep}{0in}

% Sections formatting
\titleformat{\section}{
  \vspace{-4pt}\scshape\raggedright\large
}{}{0em}{}[\color{black}\titlerule \vspace{-5pt}]

% Ensure that generate pdf is machine readable/ATS parsable
\pdfgentounicode=1

%-------------------------
% Custom commands
\newcommand{\resumeItem}[1]{
  \item\small{
    {#1 \vspace{-2pt}}
  }
}

\newcommand{\resumeSubheading}[4]{
  \vspace{-2pt}\item
    \begin{tabular*}{1.002\textwidth}[t]{l@{\extracolsep{\fill}}r}
      \textbf{#1} & #2 \\
      \textit{\small#3} & \textit{\small #4} \\
    \end{tabular*}\vspace{-7pt}
}

\newcommand{\resumeSubSubheading}[2]{
    \item
    \begin{tabular*}{0.97\textwidth}{l@{\extracolsep{\fill}}r}
      \textit{\small#1} & \textit{\small #2} \\
    \end{tabular*}\vspace{-7pt}
}

\newcommand{\resumeProjectHeading}[2]{
    \item
    \begin{tabular*}{0.99\textwidth}{l@{\extracolsep{\fill}}r}
      \small#1 & \textit{\small #2} \\
    \end{tabular*}\vspace{-7pt}
}

% NEW: Projects line format: "Project Name | Org" on left, Date on right
\newcommand{\resumeProjectSubheadingInline}[3]{
  \vspace{-0.5pt}\item
    \begin{tabular*}{1.002\textwidth}[t]{l@{\extracolsep{\fill}}r}
      \textbf{#1} \,|\, \textit{\small #2} &{\small #3} \\
    \end{tabular*}\vspace{-8pt}
}

\newcommand{\resumeSubItem}[1]{\resumeItem{#1}\vspace{-4pt}}

\renewcommand\labelitemii{$\vcenter{\hbox{\tiny$\bullet$}}$}

\newcommand{\resumeSubHeadingListStart}{\begin{itemize}[leftmargin=0in, label={}]}
\newcommand{\resumeSubHeadingListEnd}{\end{itemize}}
\newcommand{\resumeItemListStart}{\begin{itemize}[leftmargin=0.15in]}
\newcommand{\resumeItemListEnd}{\end{itemize}\vspace{-12pt}}

\newcommand{\resumeEducationCourses}[3]{
  \vspace{-2pt}\item
    \begin{tabular*}{1.002\textwidth}[t]{l@{\extracolsep{\fill}}r}
      \textbf{#1} & #2 \\
    \end{tabular*}\vspace{2pt} % <-- Controls space between title line & relevant courses
    #3\vspace{-3pt}
}

%-------------------------------------------
%%%%%%  CV STARTS HERE  %%%%%%%%%%%%%%%%%%%%%%%%%%%%


\begin{document}

%----------HEADING-----------------
\begin{center}
    \textbf{\Huge \scshape Pranay Oza} \\ \vspace{1pt}
    \small 346-630-1475 $|$ \href{mailto:oza.pranay06@gmail.com}{\underline{oza.pranay06@gmail.com}} $|$
    \href{https://www.linkedin.com/in/pranayoza/}{\underline{linkedin.com/in/pranayoza/}} $|$
    \href{https://github.com/pranay-o}{\underline{github.com/pranay-o}} $|$
    \href{https://pranay-o.github.io/}{\underline{Portfolio}}
\end{center}
\vspace{-22pt}

%-----------SKILLS-----------
\section{Skills}
\vspace{0pt}
 \begin{itemize}[leftmargin=0in, label={}]
    \small{\item{
     \textbf{Programming \& Software}{: C, C++, Assembly, SystemVerilog, Rust, Python, Java, MATLAB, Git, Unix, CMake} \\
     \textbf{Embedded Firmware}{: STM32, FreeRTOS, FDCAN/CAN, SPI/isoSPI, I\textsuperscript{2}C, UART, USB, OCPP, J-Link} \\
     \textbf{Hardware \& Systems Design}{: Altium, KiCad, Power Electronics, Circuit Analysis, Simulink, LTspice} \\
    }
    }
 \end{itemize}
 \vspace{-17pt}

%-----------EDUCATION-----------
\section{Education}
    \vspace{-5pt}
  \resumeSubHeadingListStart
    \resumeSubheading
      {University of British Columbia}{Expected Graduation -- Dec 2028}
      {Computer Engineering}{Vancouver, BC}
  \resumeSubHeadingListEnd
  \vspace{-10pt}

%-----------WORK EXPERIENCE-----------
\section{Experience}
    \vspace{-5pt}
  \resumeSubHeadingListStart

\resumeSubheading
      {High Voltage Firmware Integration}{Incoming May 2027}
      {Tesla}{Palo Alto, CA}
      \vspace{-15pt}
\resumeSubheading
      {Battery Management System Firmware Lead}{Sept 2024 -- Present}
      {UBC Formula Electric SAE}{Vancouver, BC}
      \resumeItemListStart
        \resumeItem{Leading firmware development for the car's fully custom \textbf{600V BMS}, owning the full stack from high-level \textbf{battery algorithms} down to low-level drivers monitoring 500+ 18650 lithium-ion cells}

        \resumeItem{Developing a driver for the \textbf{ADBMS2950} to monitor tractive-system voltage and pack current across a \textbf{shunt resistor}, providing the isolated high-voltage measurements required for accurate power and energy tracking}

        \resumeItem{Modeling and validating a \textbf{second-order RC} equivalent-circuit battery model in \textbf{MATLAB}, parameterized from \textbf{HPPC}, \textbf{OCV--SOC}, and capacity characterization tests ran on a battery cycler, improving fidelity of state estimation}

        \resumeItem{Researching battery state-estimation algorithms spanning \textbf{state-of-charge}, \textbf{state-of-health}, \textbf{state-of-energy}, and \textbf{state-of-power}, giving greater visibility into pack condition and to inform energy strategy}

        \resumeItem{Building \textbf{Rust} tooling that auto-generates \textbf{multiplexed FDCAN signals} from a single JSON definition and integrates with the telemetry backend, packing more data into FDCAN messages to raise logging density without adding bus load}

        \resumeItem{Relocating \textbf{ISR} and the \textbf{interrupt vector table} into \textbf{Tightly-Coupled Memory}, minimizing instruction-fetch latency to achieve deterministic, lowest-possible interrupt response times}

        \resumeItem{Validated safety-critical fault handling by injecting faults into the BMS and confirming each tripped the pack's \textbf{shutdown loop}, covering insulation-monitoring, brake-system, and inter-segment faults to ensure the pack fails safe}

        \resumeItem{Developed \textbf{CAN-based charging firmware} for a 6.6 kW Elcon charger with a full CC/CV profile and termination logic, and integrated \textbf{J1772 EVSE} charging by decoding the control-pilot duty cycle to safely limit the current request}

        \resumeItem{Implemented a \textbf{FDCAN-based bootloader} with \textbf{CRC32} verified-boot to flash all boards over FDCAN without a debugger, using \textbf{flash-memory partitioning} to keep every node on a version-matched build}

        \resumeItem{Built an \textbf{automated calibration script} and lab-bench setup for the pack current sensor, delivering the accurate power measurement the vehicle-dynamics team relied on to set the car's power limit}

        \resumeItem{Owned bring-up and on-hardware validation of the firmware on the team's custom BMS PCBs, verifying reliability through hardware testing and \textbf{software-in-the-loop (SIL)} simulation before vehicle integration}
      \resumeItemListEnd
  \resumeSubHeadingListEnd

%-----------PROJECTS-----------
\section{Projects}
\vspace{-3pt}
\resumeSubHeadingListStart
    \resumeProjectHeading
      {\normalsize\textbf{ADBMS6830B Driver Development} -- \small\textit{UBC Formula Electric SAE}}{\textup{Jan 2026 -- Sept 2026}}
      \resumeItemListStart
        \resumeItem{Owned the architecture, implementation, and integration of a driver for daisy-chained \textbf{ADBMS6830B} monitor ICs (SPI/isoSPI) across a multiple \textbf{56 cell segments} with \textbf{PEC15/PEC10 CRC} frame validation}
        \resumeItem{Implemented \textbf{cell-voltage} and \textbf{thermistor} sensing, \textbf{open-wire} detection, and per-cell \textbf{balancing} with \textbf{reversible isoSPI} and dropped-command detection for daisy-chain fault tolerance, organized into reusable modules behind a \textbf{hardware-abstraction layer} that broadcasts per-segment fault flags over FDCAN}
        \resumeItem{Designed a \textbf{FreeRTOS} concurrency model with independently clocked config, voltage, open-wire, and thermistor tasks sharing one SPI bus through \textbf{synchronization primitives}, eliminating bus contention and data races}
        \resumeItem{Debugged \textbf{15+ boards} and daisy chain for signal-integrity, hardware shorts, and DMA/SPI/task-scheduling bugs using SEGGER \textbf{SystemView} for RTOS task visibility, then ran full-battery integration testing}
      \resumeItemListEnd
      \vspace{5pt}


  \resumeProjectHeading
    {\normalsize\textbf{Sensorless Field-Oriented Control ESC}}{\textup{Jun 2025 -- Present}}
    \resumeItemListStart
        \resumeItem{Designing a \textbf{4-layer PCB} integrating an STM32G4 MCU, gate driver, buck converter, EEPROM, ESD protection circuitry, and reverse polarity detection for 30A, 12V BLDC motor control}
        \resumeItem{Developing motor-control firmware (PlatformIO, STM32CubeMX) using ADC sampling, PWM timers, I\textsuperscript{2}C, and UART for current sensing, \textbf{virtual neutral-voltage} measurement, and \textbf{back-EMF} zero-cross detection}
        \resumeItem{Building a Simulink inverter model to simulate 3-phase dynamics and validate \textbf{sensorless FOC} algorithms, including SVPWM and dead-time compensation}
      \resumeItemListEnd
    \vspace{-7pt}

\resumeSubHeadingListEnd
\vspace{2pt}
\end{document}
